\chapter{Quant Basics for Developers}
\label{chap:quantbasics}

See~\cite{hull-options},~\cite{shreve-finance},~\cite{wilmott-quant}.
\newline
This chapter assumes you know \highlight{math and programming}, but \emph{not} finance. 
It builds ideas in order --- like a physics course where each section uses only what came 
before. Formulas are there for precision; the prose is what you should read first.

\subsection*{How to read this chapter}
\begin{enumerate}
    \item \textbf{Sections~\ref{sec:plain-options}--\ref{sec:plain-volatility}:} 
          What is being traded and what ``price an option'' means.
    \item \textbf{Sections~\ref{sec:sde-intuition}--\ref{sec:bs-pde}:} 
          Why the stock is modelled as a diffusion, and how hedging leads to the 
          Black--Scholes equation (the heart of the theory).
    \item \textbf{Sections~\ref{sec:bs-closed-form}--\ref{sec:greeks}:} 
          The closed-form price and the \emph{greeks} (sensitivities) --- 
          partial derivatives, exactly like in perturbation theory.
    \item \textbf{Sections~\ref{sec:pnl-explain} onward:} 
          How desks use this in practice (P\&L, smiles, code).
\end{enumerate}
Each section has a stable \texttt{\textbackslash label} so it can be expanded later 
into full textbook depth without moving content.

\section{Stocks and Options in Plain Language}
\label{sec:plain-options}

\subsection{Stock}
A \highlight{stock} (share) is a piece of ownership in a company. Its \highlight{spot price} 
$S$ is what the market pays \emph{right now} for one share. $S$ moves over time --- that 
randomness is what makes derivatives non-trivial.

\subsection{Derivative / option}
A \highlight{derivative} is a contract whose \highlight{payoff} depends on something else 
(the \highlight{underlying}, usually the stock). The simplest running example is a 
\highlight{European call option}:
\begin{itemize}
    \item Today you agree on a \highlight{strike} $K$ and an \highlight{expiry} date $T$.
    \item At $T$, if the stock is above $K$, the call pays $(S_T - K)$; otherwise it pays $0$.
    \item In symbols: payoff $= \max(S_T - K,\, 0)$.
\end{itemize}
A \highlight{European put} pays $\max(K - S_T,\, 0)$ --- you gain if the stock \emph{falls} 
below $K$. ``European'' means you only settle at $T$, not before.

\subsection{Why anyone buys an option}
\begin{itemize}
    \item \textbf{Insurance:} a put limits losses if you own stock and the market crashes.
    \item \textbf{Leverage:} a call gives exposure to upside with limited upfront cost 
          (the \highlight{premium}, i.e.\ the option's price today).
    \item \textbf{Trading / hedging:} desks replicate payoffs, trade volatility, or hedge 
          other books --- your job as quant dev is to \highlight{price} and \highlight{risk-manage} 
          these contracts in code.
\end{itemize}
\textbf{Key idea:} the option has a \highlight{value today}, $V(S,t)$, even though the 
payoff is uncertain. Pricing = finding a fair $V$ given $S$, $K$, $T$, and model assumptions.

\section{Interest Rates, Dividends, and Forward Price}
\label{sec:plain-rates-divs}
\begin{itemize}
    \item \highlight{$r$} --- \textbf{risk-free rate} (annualized). Money in the bank grows 
          as $e^{rt}$; discounting future cash by $e^{-r(T-t)}$ is the continuous analogue 
          of ``\$1 today vs \$1 in six months.''
    \item \highlight{$q$} --- \textbf{dividend yield}. Shareholders receive dividends; 
          holding stock is slightly less attractive than holding cash at $r$, so the 
          \emph{effective} drift of $S$ in no-arbitrage pricing uses $(r - q)$, not $r$ alone.
    \item \textbf{Forward price} (simplified): 
          $F = S\,e^{(r-q)(T-t)}$ is the fair locked-in price to deliver stock at $T$ 
          agreed today. Many desks think in forwards rather than spots.
\end{itemize}
\par\noindent\textbf{Interview tip:} you do not need fixed-income expertise for a first pass --- just know that 
$e^{-rT}$ discounts and $q$ shifts the fair drift.

\section{Volatility}
\label{sec:plain-volatility}
\highlight{Volatility} $\sigma$ measures how much $S$ \emph{fluctuates}. In the models below, 
$\sigma$ is the coefficient on the random noise (like a diffusion constant). 
It is quoted as a \textbf{decimal}: $\sigma = 0.20$ means ``20\% per year'' in the 
standard annualized sense.
\newline
Higher $\sigma$ $\Rightarrow$ more uncertainty $\Rightarrow$ an option (which pays off 
in the tails) is usually \textbf{worth more}. That is the economic content of \highlight{vega} 
(Section~\ref{sec:greeks}).

\section{What a Quant Developer Does}
\label{sec:quant-dev-role}
At systematic trading firms, a \highlight{Senior Quant Developer} typically:
\begin{itemize}
    \item Writes C++ (or similar) to \textbf{price} options and compute \textbf{greeks}.
    \item Makes libraries \textbf{fast, deterministic, and tested} on paths that affect 
          live \highlight{P\&L} (profit and loss).
    \item Connects \textbf{market data} (prices, rates, vol surfaces) to those libraries.
    \item Explains daily P\&L: ``how much was stock move vs vol move vs time decay?''
\end{itemize}
\par\noindent\textbf{Interview tip:} understand \emph{both} the model and \emph{where production breaks} 
(wrong units, stale data, smile not flat, etc.).

\section{Modelling the Stock: Random Walk and GBM}
\label{sec:sde-intuition}
\textbf{Physics analogy:} treat $S_t$ like a stochastic process with drift and noise.

\subsection{Discrete warm-up}
Over a small time step $\Delta t$, a crude model is:
\[
  \frac{\Delta S}{S} \approx \mu\,\Delta t + \sigma\,\sqrt{\Delta t}\,\xi,
  \qquad \xi \sim \mathcal{N}(0,1).
\]
$\mu$ is average relative return; $\sigma$ controls the size of shocks. 
Taking $\Delta t \to 0$ gives \highlight{geometric Brownian motion} (GBM):
\begin{equation}
  dS_t = \mu S_t\,dt + \sigma S_t\,dW_t ,
  \label{eq:gbm}
\end{equation}
with $W_t$ standard Brownian motion. Because the noise multiplies $S_t$, $S$ stays 
\highlight{positive} and is \highlight{log-normal} at fixed time (not Gaussian).

\subsection{Risk-neutral measure (pricing only)}
For \textbf{pricing}, we do not use your real-world $\mu$ (which is hard to estimate). 
No-arbitrage implies we can price as if, under a special probability measure 
$\mathbb{Q}$,
\begin{equation}
  dS_t = (r - q)\,S_t\,dt + \sigma S_t\,dW_t^{\mathbb{Q}} .
  \label{eq:gbm-rn}
\end{equation}
\textbf{In words:} in the pricing world, the \emph{expected} growth rate of $S$ 
(adjusted for dividends) equals the risk-free rate. The discounted option value 
$e^{-rt}V_t$ is then a \highlight{martingale} --- expected future payoff, 
discounted, equals today's price. This is the finance version of ``choose coordinates 
so the equation simplifies.''

\section{Black--Scholes: Assumptions and the Big Idea}
\label{sec:bs-assumptions}
Black--Scholes (BS)~\cite{wiki-black-scholes} is the baseline model. It assumes:
\begin{itemize}
    \item $S$ follows GBM with \textbf{constant} $\sigma$.
    \item You can trade $S$ continuously, no fees, borrow/lend at rate $r$.
    \item No jumps, no smile built in (vol same for all strikes).
\end{itemize}
\textbf{The big idea (replication):} if you hold the option and continuously adjust a 
position in the stock, you can \highlight{cancel the leading risk} of $dS$. 
The remaining portfolio must earn the risk-free rate (else arbitrage). 
That balance \emph{forces} a PDE on $V(S,t)$ --- you are not guessing a price; 
you are solving for the unique price consistent with no arbitrage + hedging.

\section{The Black--Scholes PDE}
\label{sec:bs-pde}
Let $V(S,t)$ be the option value. \textbf{It\^o's lemma} (chain rule with $dW$) gives 
how $dV$ depends on $dS$ and $dt$. A portfolio ``long 1 option, short $\Delta$ shares'' 
can be made insensitive to $dS$ to first order by choosing $\Delta = \partial V/\partial S$. 
No-arbitrage then yields:
\begin{equation}
  \frac{\partial V}{\partial t}
  + (r - q)\,S\,\frac{\partial V}{\partial S}
  + \tfrac{1}{2}\sigma^2 S^2\,\frac{\partial^2 V}{\partial S^2}
  - r\,V = 0 .
  \label{eq:bs-pde}
\end{equation}
\textbf{Read it like a physics PDE:} time derivative + drift in $S$ + diffusion in $S$ 
($\tfrac{1}{2}\sigma^2 S^2 \partial_{SS}$) + discount term $-rV = 0$. 
\textbf{Terminal condition} at $T$: $V(S,T) = \max(S-K,0)$ for a call. 
Solve backward in time from $T$ to today.

\section{Closed-Form Price (European Call and Put)}
\label{sec:bs-closed-form}
BS gives an explicit solution (no need to solve the PDE numerically for Europeans). Define:
\begin{align}
  d_1 &= \frac{\ln(S/K) + (r - q + \tfrac{1}{2}\sigma^2)(T-t)}{\sigma\sqrt{T-t}} ,
  &
  d_2 &= d_1 - \sigma\sqrt{T-t} .
  \label{eq:d1d2}
\end{align}
$N(x)$ = standard normal CDF, $\phi(x)$ = PDF. Then:
\begin{align}
  C &= S\,e^{-q(T-t)}\,N(d_1) - K\,e^{-r(T-t)}\,N(d_2) \quad\text{(call)}, \\
  P &= K\,e^{-r(T-t)}\,N(-d_2) - S\,e^{-q(T-t)}\,N(-d_1) \quad\text{(put)}.
\end{align}
\textbf{Rough intuition:}
\begin{itemize}
    \item $N(d_2)$ is (roughly) the risk-neutral probability the call finishes in the money.
    \item $N(d_1)$ weights the stock leg of the replicating portfolio (the $\Delta$ hedge).
    \item $e^{-r(T-t)}$ and $e^{-q(T-t)}$ discount strike and spot contributions.
\end{itemize}
\par\noindent\textbf{Interview tip:} you do not need to derive $d_1,d_2$ --- know they encode 
\highlight{moneyness} ($S/K$), \highlight{time to expiry}, \highlight{rates}, and \highlight{vol}.

\subsection{Put--Call Parity}
\label{subsec:put-call-parity}
For the same $K$, $T$~\cite{wiki-put-call-parity}:
\begin{equation}
  C - P = S\,e^{-q(T-t)} - K\,e^{-r(T-t)} .
  \label{eq:put-call-parity}
\end{equation}
\textbf{In words:} ``long call + short put'' replicates a forward on the stock. 
If this fails in your pricer, you have a bug or inconsistent inputs.

\section{The Greeks: Sensitivities of the Price}
\label{sec:greeks}
\highlight{Greeks}~\cite{wiki-greeks} are partial derivatives of $V$ --- exactly like 
sensitivity of an observable to parameters in physics.

\begin{center}
\begin{tabular}{|l|l|p{7.5cm}|}
\hline
\textbf{Name} & \textbf{Symbol} & \textbf{Meaning in one sentence} \\
\hline
Delta   & $\Delta = \partial V/\partial S$ 
        & How much option value moves when $S$ moves \$1 (hedge ratio in shares). \\
Gamma   & $\Gamma = \partial^2 V/\partial S^2$ 
        & How fast $\Delta$ changes when $S$ moves (curvature / convexity). \\
Vega    & $\nu = \partial V/\partial \sigma$ 
        & Sensitivity to vol; long options usually have $\nu > 0$. \\
Theta   & $\Theta = \partial V/\partial t$ 
        & Time decay; long vanilla usually loses value as $T$ approaches. \\
Rho     & $\rho = \partial V/\partial r$ 
        & Sensitivity to interest rates (often smaller for short-dated equity options). \\
\hline
\end{tabular}
\end{center}
For a European \textbf{call} under BS:
\begin{align}
  \Delta &= e^{-q(T-t)}\,N(d_1) , &
  \Gamma &= \frac{e^{-q(T-t)}\,\phi(d_1)}{S\,\sigma\sqrt{T-t}} , \\
  \nu    &= S\,e^{-q(T-t)}\,\phi(d_1)\,\sqrt{T-t} .
\end{align}
\textbf{Units trap:} $\sigma = 0.20$ is 20\%; ``vega per vol \emph{point}'' often means 
$\nu/100$ (per 1 percentage point).

\subsection{Delta hedging}
\label{subsec:delta-hedging}
If you are \textbf{short} one call (you sold it), you are exposed to $S$ going up. 
Hold $\Delta$ shares long: for small moves, option loss $\approx$ stock gain. 
\textbf{Delta-neutral} = net $\Delta$ of book is zero. You re-hedge when $S$ moves 
(discrete hedging) --- real P\&L differs from the idealized continuous model.

\subsection{Gamma and theta (long option)}
\label{subsec:gamma-theta}
\textbf{Long} a call: $\Gamma > 0$ (convexity --- you gain more from up moves than you 
lose from down moves, to second order), but $\Theta < 0$ (you \emph{pay} for that 
convexity via time decay). \textbf{Short} option: opposite signs --- you collect theta 
but suffer if $S$ moves a lot (gamma risk). Near expiry at-the-money, $\Gamma$ spikes 
--- hedging is hardest there.

\section{P\&L Explain: Accounting for Daily Profit and Loss}
\label{sec:pnl-explain}
Desks decompose one day's P\&L using a Taylor expansion --- same logic as 
``$\delta(\mathrm{observable}) \approx (\partial O/\partial p)\,\delta p$'':
\begin{equation}
  \Delta \mathrm{PnL} \approx
  \Delta\,\delta S
  + \tfrac{1}{2}\Gamma\,(\delta S)^2
  + \nu\,\delta\sigma
  + \Theta\,\delta t
  + \rho\,\delta r
  + \text{residual} .
  \label{eq:pnl-explain}
\end{equation}
\textbf{Concrete reading:}
\begin{itemize}
    \item $\Delta\,\delta S$ --- P\&L from the stock moving (delta explain).
    \item $\tfrac{1}{2}\Gamma\,(\delta S)^2$ --- correction when the move was large 
          (gamma explain).
    \item $\nu\,\delta\sigma$ --- P\&L from implied vol changing (vol explain).
    \item $\Theta\,\delta t$ --- P\&L from one day passing (time decay).
    \item \textbf{Residual} --- anything left: bad data, model mismatch, cross terms, 
          discrete hedging, dividends not in model, etc.
\end{itemize}
\par\noindent\textbf{Interview tip:} ``We lost \$2M; \$1.6M is delta on a $-3\%$ move, 
\$300k gamma, \$100k vega when skew steepened; \$50k unexplained --- we should check 
surface build.''

\section{Implied Volatility and the Smile}
\label{sec:implied-vol-smile}
BS uses one number $\sigma$ for all strikes. \textbf{Markets do not:} if you take 
each traded option price and ask ``what $\sigma$ makes BS match?'', you get 
\highlight{implied volatility} $\sigma_{\mathrm{imp}}(K,T)$ --- typically a 
\highlight{smile/skew}~\cite{wiki-volatility-smile} (equity puts often have higher 
implied vol than calls at same expiry).
\newline
\textbf{In plain terms:} BS is a wrong model, but $\sigma_{\mathrm{imp}}$ is a 
\highlight{useful coordinate} for quoting options (like using action angle variables 
when the true Hamiltonian is messy). Production stores a \highlight{vol surface}, 
not one scalar $\sigma$.

\section{Beyond Vanilla (Short Glossary)}
\label{sec:beyond-vanilla-glossary}
Read these once; depth can come later.
\begin{itemize}
    \item \textbf{Realized vol} --- measured from past returns; \textbf{implied vol} --- 
          inferred from option prices.
    \item \textbf{Exotic} --- non-vanilla payoff (barrier, Asian average, digital); 
          often priced by Monte Carlo or PDE on a grid.
    \item \textbf{American} --- can exercise before $T$; worth at least European.
    \item \textbf{Vanna / volga} --- cross and second derivatives w.r.t.\ $S$ and $\sigma$; 
          matter on large joint moves.
    \item \textbf{Bid--ask / order book} --- where theoretical price meets tradable price 
          (Part~II networking chapters).
\end{itemize}

\section{Computing Greeks in Code}
\label{sec:analytic-vs-bump}
Two approaches:
\begin{itemize}
    \item \textbf{Analytic:} plug into closed forms above (fast for BS).
    \item \textbf{Bump:} reprice with $S \to S+h$, $\sigma \to \sigma+\varepsilon$; 
          finite difference (works for any pricer):
\end{itemize}
\begin{align}
  \Delta &\approx \frac{V(S+h) - V(S-h)}{2h} , &
  \nu    &\approx \frac{V(\sigma+\varepsilon) - V(\sigma-\varepsilon)}{2\varepsilon} .
\end{align}
Choose $h$ like any numerical derivative: not too small (roundoff) or too large (truncation).

\toclessnparagraph{Minimal C++ pricer sketch}
\begin{lstlisting}[language=C++]
#include <cmath>

struct BsParams { double S, K, r, q, sigma, T; };

inline double norm_cdf(double x) {
    return 0.5 * std::erfc(-x / std::sqrt(2.0));
}

double bs_call(const BsParams& p) {
    const double sqT = std::sqrt(p.T);
    const double d1 = (std::log(p.S / p.K) + (p.r - p.q + 0.5 * p.sigma * p.sigma) * p.T)
                      / (p.sigma * sqT);
    const double d2 = d1 - p.sigma * sqT;
    return p.S * std::exp(-p.q * p.T) * norm_cdf(d1)
         - p.K * std::exp(-p.r * p.T) * norm_cdf(d2);
}

double bump_delta(const BsParams& p, double h = 1e-4 * p.S) {
    BsParams up = p; up.S += h;
    BsParams dn = p; dn.S -= h;
    return (bs_call(up) - bs_call(dn)) / (2.0 * h);
}
\end{lstlisting}
Use \texttt{std::erfc}~\cite{cppreference-erfc} for stable $N(x)$ in the tails.

\section{Implementation Pitfalls}
\label{sec:numerical-gotchas}
\begin{itemize}
    \item \textbf{Near expiry ATM:} greeks blow up; use limits or switch to intrinsic value.
    \item \textbf{Units:} vol decimal vs percent; rates; day-count conventions.
    \item \textbf{Put--call parity:} automated test on every build.
    \item \textbf{Determinism:} Monte Carlo greeks need fixed seeds for regression tests.
\end{itemize}

\section{Monte Carlo and PDE (When BS Is Not Enough)}
\label{sec:monte-carlo}
\textbf{Monte Carlo}~\cite{wiki-monte-carlo}: simulate many paths of $S$, average discounted 
payoffs --- works for path-dependent exotics; error $\propto 1/\sqrt{N}$.
\newline
\textbf{PDE / finite differences:} discretize Eq.~\eqref{eq:bs-pde} on a grid in $(S,t)$; 
needed for Americans and some barriers. As quant dev you choose method + validate against 
known cases (BS, parity).

\section{Production Architecture (Lite)}
\label{sec:pricer-risk-design}
Typical pipeline: \textbf{market data snapshot} $\to$ \textbf{curves/surfaces} $\to$ 
\textbf{pricer} $\to$ \textbf{risk aggregator}. Version snapshots so you can replay 
``what did we think P\&L was at 10:00?'' See Chapter~\ref{chap:lowlatency} for hot-path 
engineering and Chapter~\ref{chap:concurrency} for thread-safe shared state.

\section{Testing and Python/C++ Split}
\label{sec:testing-quant-code}
Test: parity, analytic vs bump greeks, edge cases ($T \to 0$, deep OTM). 
Research in Python; production pricers in C++ with a thin binding --- same numerics, 
stricter latency and determinism requirements.

\section{Interview Cheat Sheet (Read Last)}
\label{sec:quant-interview-cheat-sheet}
After reading the chapter once, rehearse these aloud:
\begin{itemize}
    \item \textbf{Option:} right to receive $\max(S_T-K,0)$ at $T$; priced by no-arbitrage 
          + hedging, not by guessing $\mu$.
    \item \textbf{BS:} GBM, constant $\sigma$, replicate with $\Delta$ shares $\Rightarrow$ PDE 
          $\Rightarrow$ closed form with $N(d_1), N(d_2)$.
    \item \textbf{Greeks:} $\partial V$ w.r.t.\ $S,\sigma,t,r$; delta hedge; long option = 
          +gamma, $-$theta usually.
    \item \textbf{Smile:} one $\sigma$ is wrong; desks use surfaces; implied vol is a quote 
          coordinate.
    \item \textbf{P\&L explain:} Taylor in $\delta S$, $\delta\sigma$, $\delta t$; residual 
          = bug or model gap.
    \item \textbf{Dev:} test parity, mind units, bump vs analytic, deterministic MC.
\end{itemize}

\section{Typical Interview Questions}
\label{sec:quant-interview-qa}
Short answers you can expand on a whiteboard. For order-book data structures in more depth, 
see also Chapter~\ref{chap:algorithms} (order book section).

\subsection{Market data and the order book}

\toclessnparagraph{What is the difference between Level~1, Level~2, and Level~3 market data?}
\begin{itemize}
    \item \highlightbf{Level~1 (L1) / top of book:} best available prices only. Typically 
          best bid price/size, best ask price/size, and often last trade price/volume. 
          You see the \highlight{spread} and the last print, but \emph{not} book depth.
    \item \highlightbf{Level~2 (L2) / market depth:} many price levels with 
          \highlight{aggregated} size at each price (not every order id). Example:
\begin{center}
\small
\begin{tabular}{lr@{\hspace{1.5em}}lr}
\multicolumn{2}{c}{\textbf{Ask}} & \multicolumn{2}{c}{\textbf{Bid}} \\
101.30 & 20 & 101.27 & 15 \\
101.29 & 35 & 101.26 & 40 \\
101.28 & 50 & 101.25 & 60 \\
\end{tabular}
\end{center}
          From L2 you can read depth, liquidity, bid/ask \highlight{imbalance}, and 
          short-horizon pressure. Many HFT strategies need at least L2; queue-position 
          strategies may need L3.
    \item \highlightbf{Level~3 (L3) / order-by-order:} individual orders with ids. 
          Instead of \texttt{101.25 $\rightarrow$ 300}, you see order 1452 buy 100, 
          1891 buy 50, 2110 buy 150, and events \texttt{New}/\texttt{Modify}/\texttt{Cancel}/\texttt{Execution}. 
          Full \highlight{limit order book} reconstruction and price--time queues are possible; 
          heaviest bandwidth and bookkeeping.
\end{itemize}
Naming varies by venue/vendor; say what \emph{your} feed actually contains.

\toclessnparagraph{What information is contained in a market data feed?}
Depends on the feed. Common pieces:
\begin{itemize}
    \item \textbf{Order events:} new / modify / cancel / replace (L3) or level updates (L2).
    \item \textbf{Trade events:} execution price and size.
    \item \textbf{Quotes:} resting interest or aggregated BBO (an intention, not a fill).
    \item \textbf{Market state:} halt, auction, open, close.
    \item \textbf{Timestamps:} exchange / matching / send --- critical for HFT 
          (often µs or ns resolution, venue-dependent).
    \item \textbf{Sequence numbers:} gap detection and replay.
\end{itemize}
Options/futures feeds may add extra fields (e.g.\ open interest) --- venue-dependent.

\toclessnparagraph{What is the difference between trades and quotes?}
A \highlight{quote} is an \emph{intention}: e.g.\ buy 100 @ 100.20 or sell 200 @ 100.25 
--- resting (or displayed) interest, not necessarily a fill. It may arrive as full 
orders (L3) or as aggregated BBO/level updates.  
A \highlight{trade} is an \emph{execution}: e.g.\ 50 @ 100.23 --- someone bought and 
someone sold.  
\textbf{Short:} quotes $\approx$ interest in the book (or top-of-book updates); 
trades $\approx$ prints. Strategies use quotes for shape and trades for markouts.

\toclessnparagraph{How would you reconstruct an order book from incremental updates?}
Modern feeds rarely blast the full book every time; they send an \highlight{event stream}.
\begin{enumerate}
    \item Start from an initial \highlight{snapshot} (or empty book if the protocol allows).
    \item Apply sequenced messages locally: \texttt{Add} (add liquidity), \texttt{Modify} 
          (size/price), \texttt{Cancel} (remove), \texttt{Execution}/\texttt{Trade} 
          (reduce or delete the resting order that was hit).
    \item Keep invariants: no empty levels; BBO consistent; detect \highlight{sequence gaps} 
          and resubscribe/snapshot.
\end{enumerate}
Incremental updates are far cheaper than continuous full snapshots. Never assume 
cross-feed wall-clock order without sequence checks.

\toclessnparagraph{Why would an HFT strategy need Level~2 instead of Level~1?}
L1 only shows the touch. Two books can share the \highlight{same best ask} and look identical 
on L1 while liquidity behind the touch is completely different:
\begin{center}
\small
\begin{tabular}{lr@{\hspace{2em}}lr}
\multicolumn{2}{c}{\textbf{Ask A}} & \multicolumn{2}{c}{\textbf{Ask B}} \\
101.00 & 5   & 101.00 & 500 \\
101.01 & 800 & 101.01 & 5 \\
101.02 & 900 & 101.02 & 5 \\
\end{tabular}
\end{center}
L2 unlocks depth, how far a take walks the book, and signals such as 
\highlight{order-book imbalance}, liquidity, pressure, support/resistance style levels, 
and \highlight{microprice}. Many quantitative short-horizon strategies use exactly these.

\toclessnparagraph{How would you represent an order book in memory?}
State the hot query first (BBO vs deep walk vs cancel-by-id). A common L3 sketch:
\begin{lstlisting}[language=C++]
// bids: highest price first; asks: lowest first
std::map<Price, PriceLevel, std::greater<>> bids;
std::map<Price, PriceLevel>                 asks;
// PriceLevel: price, total volume, queue of orders (price-time)
std::unordered_map<OrderId, Order*>         by_id;  // O(1) modify/cancel
\end{lstlisting}
Each price level holds the resting orders at that price; \texttt{by\char`_id} makes 
modify/cancel fast.  
In real HFT stacks, generic \texttt{std::map} is often avoided when the tick grid or 
price range is bounded: prefer \highlight{cache-friendly} layouts (tick-indexed arrays, 
custom trees, pooled nodes) to cut allocations and improve locality. 
See Chapter~\ref{chap:algorithms}.

\toclessnparagraph{What is price--time priority?}
At the same price, earlier resting orders trade first. Your L3 book must preserve 
queue order if you model fill probability at the touch.

\toclessnparagraph{Market order vs limit order?}
\highlight{Market:} take liquidity now at whatever prices are available (walk the book). 
\highlight{Limit:} rest at a price (or take if marketable). HFT cares about 
\highlight{make} vs \highlight{take}, fees, and queue position.

\subsection{Latency and systems (quant-dev)}

\toclessnparagraph{How is tick-to-trade latency composed?}
\label{qa:latency-composition}
Always define the path first. The industry default is \highlight{tick-to-trade}:
\textbf{market-data packet in} $\rightarrow$ \textbf{order packet out} toward the venue.
\newline
\begin{center}
% Tick-to-trade latency pipeline (externalized)
\centering
\tikzfig{tick_to_trade_latency}
\begin{tikzpicture}[
    font=\sffamily\footnotesize,
    stage/.style={
        draw=black,
        thick,
        rounded corners=1pt,
        align=center,
        minimum height=1.15cm,
        text width=1.55cm,
        inner sep=2pt
    },
    arr/.style={-{Stealth[length=2mm]}, thick},
    note/.style={font=\sffamily\scriptsize, align=center, text=oblivioncomment}
]
    \node[stage] (nic) {NIC /\\kernel};
    \node[stage, right=0.35cm of nic] (dec) {Decode /\\normalize};
    \node[stage, right=0.35cm of dec] (book) {Book\\update};
    \node[stage, right=0.35cm of book] (strat) {Strategy};
    \node[stage, right=0.35cm of strat] (risk) {Risk\\checks};
    \node[stage, right=0.35cm of risk] (enc) {Encode\\+ send};

    \draw[arr] (nic) -- (dec);
    \draw[arr] (dec) -- (book);
    \draw[arr] (book) -- (strat);
    \draw[arr] (strat) -- (risk);
    \draw[arr] (risk) -- (enc);

    \node[note, above=0.15cm of nic] {MD in};
    \node[note, above=0.15cm of enc] {order out};

    % harness overlay note
    \node[note, below=0.55cm of book] (harness) {%
        \textbf{Harness variant:} inject simulated MD here $\rightarrow$ stop at simulated order\\
        (app path only --- not exchange RTT)};
    \draw[arr, densely dashed, oblivioncomment] ([yshift=-0.05cm]dec.south) -- ++(0,-0.25)
        -| ([xshift=-0.1cm]harness.north);
\end{tikzpicture}

\captionof{figure}{Software tick-to-trade path (and harness variant).}
\label{fig:tick-to-trade-latency}
\end{center}

\textbf{Budget buckets} (whiteboard version):
\begin{center}
\small
\begin{tabular}{|l|p{8.2cm}|}
\hline
\textbf{Stage} & \textbf{What eats time} \\
\hline
NIC / kernel & Interrupt or poll, sk\_buff, copies, syscall; shrinks with busy-poll / bypass 
  (Chapters~\ref{chap:eventdrivenio},~\ref{chap:kernelbypass}). \\
\hline
Decode / normalize & Parse wire format, instrument map, sequence checks. \\
\hline
Book update & Apply L2/L3 deltas; map/array walks; allocator traffic. \\
\hline
Strategy & Signal + decision; can dominate if heavy or blocked on locks. \\
\hline
Risk / pre-trade & Limits, fat-finger, position; must stay allocation-thin. \\
\hline
Encode + send & Serialize order, \texttt{send}/\texttt{write}; Nagle off on hot TCP. \\
\hline
Cross-cutting tax & Queues, context switches, cache misses / false sharing, NUMA remote, 
  page faults, logging --- often the \emph{tail}, not the median. \\
\hline
\end{tabular}
\end{center}

\textbf{What usually dominates?}
\begin{itemize}
    \item \textbf{Fat app:} book + strategy + risk + hidden allocations.
    \item \textbf{Lean app:} kernel path, copies, cache/NUMA, NIC --- optimize the wire side 
          (Part~II) once the software path is already thin.
\end{itemize}
Interview line: give a number only with a path (``local process path $\sim$1--10\,µs is 
plausible; exchange RTT is another story'').

\toclessnparagraph{What about inject $\rightarrow$ simulated order (harness latency)?}
Lab metrics often start at \highlight{controlled injection} of market data and stop when the 
stack emits a \highlight{simulated order} (no venue ACK). That measures the 
\textbf{application path} under load and is valid if you \emph{name} it --- it is 
\emph{not} full tick-to-trade to the exchange.  
Use the same harness before/after each optimization (pinning, NUMA split, lock-free queues, 
lean serialization, monotonic clocks). Production still needs separate clocks on the 
real MD-in / order-out boundaries.

\toclessnparagraph{Median vs tail --- why HFT cares about p99?}
Median (p50) shows the happy path. Tails (p99/p999) come from interrupts, stealing, 
allocator locks, cold caches, page faults, disk logs, and noisy neighbors. A strategy can 
look fine on average and still lose on the slow prints. Optimize for 
\highlight{predictability}, not only average.

\toclessnparagraph{Where does latency come from on a tick-to-trade path?}
See the composition above (Figure~\ref{fig:tick-to-trade-latency}). Short answer: 
NIC/kernel $\rightarrow$ decode $\rightarrow$ book $\rightarrow$ strategy $\rightarrow$ 
risk $\rightarrow$ encode/send, plus jitter taxes. Book + strategy + risk often dominate 
until the app is lean.

\toclessnparagraph{How would you measure software latency?}
\begin{itemize}
    \item Define start/end events on a \highlight{named path}.
    \item Steady clock: \texttt{CLOCK\char`_MONOTONIC\char`_RAW} (or equivalent) --- not 
          \texttt{gettimeofday} for intervals.
    \item Report median \emph{and} tails; same hardware affinity before/after.
    \item Separate \textbf{harness / app path} from \textbf{exchange RTT}.
\end{itemize}

\toclessnparagraph{Why pin threads and care about NUMA?}
Keep hot threads on dedicated cores; keep memory and NIC on the \highlight{same node} 
as the consumer to avoid remote DRAM and cross-QPI/UPI traffic 
(Chapters~\ref{chap:cacheperformance},~\ref{chap:numa}).

\toclessnparagraph{How do you pass opaque payloads in C++?}
Be explicit about the layer: hot-path FFI often uses \texttt{void*} + size (or a byte 
buffer + tag). Library design may use type erasure (\texttt{std::any}, 
\texttt{std::function}, or a small virtual/SOP eraser). Ask which layer they mean.

\subsection{Pricing, risk, and greeks}

\toclessnparagraph{What is put--call parity?}
For European options with the same $K$, $T$: 
$C - P = S e^{-q(T-t)} - K e^{-r(T-t)}$. Violation $\Rightarrow$ arb or a pricer bug.

\toclessnparagraph{What is delta hedging in one sentence?}
Hold $-\Delta$ shares per short option (signs flip with position) so small moves in $S$ 
cancel to first order; rebalance as $\Delta$ changes (gamma).

\toclessnparagraph{What is implied volatility?}
The $\sigma$ that makes the BS price match the market price. One number per option; 
across strikes you get a \highlight{smile/skew}, so production uses a surface, not one 
global $\sigma$.

\toclessnparagraph{Analytic vs bump greeks?}
Analytic: closed-form $\partial V$ under BS --- fast, model-tied. 
Bump: reprice at $S\pm h$ or $\sigma\pm\varepsilon$ --- works for any pricer, careful with 
step size and noise.

\toclessnparagraph{How do you explain daily P\&L?}
Taylor: $\Delta\,\delta S + \tfrac{1}{2}\Gamma(\delta S)^2 + \nu\,\delta\sigma + \Theta\,\delta t + \cdots$ 
plus residual (data/model/hedge error).

\subsection{Behavioural / role}

\toclessnparagraph{Why move from industrial C++ (e.g.\ semis) back to trading systems?}
Tight feedback loop: ship, measure latency/correctness, iterate with the desk. 
Ownership of a hot path and measurable impact --- not month-long silicon cycles.

\toclessnparagraph{What would you ask us?}
Team size and greenfield vs reuse; what you own on the critical path; how latency is 
measured in prod; on-call expectations; what success looks like in six months.
